% GENERATED FILE. Edit canonical lessons, then run npm run generate:books.
% canonical-source-hash: fnv1a64:33fbe515044f949b
% canonical-lessons: PA-R25-drink-request-r4, PA-R25-milk-bread-r4, PA-R25-friend-family-r4, PA-R25-kin-eye-r4

\chapter{Requests, People, and Body Words at R4}
\label{ch:pa-request-people-r4-bridge}
\hlchaptermodality{25}

\begin{chapteropening}
\textbf{By the end of this chapter:} \emph{I can retrieve old request, food, people, kin, and body-word families at long distance in small spoken and reading passes.}

The learner closes the moving R4 boundary with a bounded contrast among two low-tone kin words and one doubled-consonant body word, without claiming writing evidence.
\end{chapteropening}

\section[Practice]{Water and tea requests return at R4}

\label{lesson:PA-R25-drink-request-r4}



\begin{quote}
Ask for water, then tea, using the same known polite close each time.
\end{quote}

\subsection*{Your turn}
\textbf{Two remembered routes.}

Read \textbf{\pa{ਪਾਣੀ} · \pa{ਚਾਹ}} once. This is recognition and spoken retrieval, not a writing checkpoint.

\begin{tcolorbox}[breakable,colback=teal!4,colframe=teal!35!black,title={Before you move on}]
Make the two polite requests once more from meaning.
\end{tcolorbox}
\section[Practice]{Milk and bread return at R4}

\label{lesson:PA-R25-milk-bread-r4}



\begin{quote}
Hear \textbf{dudh · roṭī} and say “milk · bread” aloud.
\end{quote}

\subsection*{Your turn}
Read \textbf{\pa{ਦੁੱਧ} · \pa{ਰੋਟੀ}} once. No copying is requested, and no independent writing evidence is awarded.

\begin{tcolorbox}[breakable,colback=teal!4,colframe=teal!35!black,title={Before you move on}]
Ask for either one using the already-known \textbf{kirpā karke}.
\end{tcolorbox}
\section[Practice]{Friend and family return at R4}

\label{lesson:PA-R25-friend-family-r4}



\begin{quote}
Say \textbf{dost}, “friend,” and \textbf{parivār}, “family.”
\end{quote}

\subsection*{Your turn}
Read \textbf{\pa{ਦੋਸਤ} · \pa{ਪਰਿਵਾਰ}} once. The models support reading only; do not count them as independent writing.

\begin{tcolorbox}[breakable,colback=teal!4,colframe=teal!35!black,title={Before you move on}]
Say “friend; family” from meaning alone.
\end{tcolorbox}
\section[Practice]{Brother, sister, and eye return at R4}

\label{lesson:PA-R25-kin-eye-r4}



\begin{quote}
Say \textbf{bharā}, \textbf{bhaiṇ}, and \textbf{akkh} from memory: brother, sister, eye. Keep the first two low; hold the doubled consonant in the third.
\end{quote}

\subsection*{Your turn}
\textbf{Family trees.}

Read \textbf{\pa{ਭਰਾ} · \pa{ਭੈਣ} · \pa{ਅੱਖ}} once. This lesson deliberately stops at reading and speech; it does not award independent Gurmukhi writing evidence.

\begin{tcolorbox}[breakable,colback=teal!4,colframe=teal!35!black,title={Before you move on}]
Say the three meanings once more, then stop. No new item follows.
\end{tcolorbox}
